% SPDX-License-Identifier: CC-BY-4.0
% Copyright (c) 2026 Martin Schröder <info@swedishembedded.com>

\section{Results}\label{sec:rsi}

\subsection{S4: retraining from the base on a doubling frontier}\label{sec:rsi4}

S4 changed what S3's records pointed at (\cref{sec:earlier}), and with it the schedule
(\cref{fig:cycle}). \emph{Retraining from the base}: generation $g$'s candidate is a fresh LoRA trained
from the base on the union of every generation's frontier, not a continuation of generation $g-1$'s
adapter; only the data of generation $g$ is new, proposed and predicted by generation $g-1$'s
champion. \emph{Frontier selection}: a program enters when at least one of the champion's 4 sampled
predictions (temperature 1.0) is wrong, so the frontier also holds programs it usually gets right: one
predicted right with probability 0.9 per sample enters with probability 0.34.\cl{RS25,RS35}
\emph{Schedule}: 300, 600 and 1200 new programs per generation, two passes over the union, no separate
replay, so the generations trained on 300, 900 and 2100 records for 75, 225 and 525 steps.\cl{RS25,RS26}
These changes were made together and no ablation separates them. Before any verdict, the examiner's
null replicate (base against base) found no discordant item, the base was served at the configuration
under which it answers 12 exam items (\cref{app:series}), and the examiner's build was changed so that
it identifies the declared base by the checkpoint's content.\cl{RS07,RS27,RS06}

\begin{table}[tbp]
\centering\small
\setlength{\tabcolsep}{3.4pt}
\input{generated/rsi_table}
\caption{S3 and S4 as the examiner decided them: items right of 889 and each step's wins/losses
with the exact one-sided sign test (a step holds at $p<0.05$ with wins above losses; the criterion is
3 of 4 lineages holding at every step). Pooled rows sum the lineage rows (they equal the plan verdict's) and are descriptive, without a test, because the lineages share the exam's items.\cl{RS30} $^\dagger$~r2 is compared on the 888 items left after one overlapping item was excluded; at generation~3 it answered 242 of the 889 before the exclusion. S4 trained on \mRsiFourDevRZeroGOneTrainedRecords{}, \mRsiFourDevRZeroGTwoTrainedRecords{} and \mRsiFourDevRZeroGThreeTrainedRecords{} cumulative records at generations 1 to 3. S3's base was served under a different configuration from S4's (\cref{app:series}), so compare steps, not levels.}\label{tab:rsi}
\end{table}

\paragraph{Outcome (adapter level).} All twelve S4 steps held (\cref{tab:rsi}). From 12 items at the
base the lineages reached 192, 225, 241 and 218 at generation~3; the weakest step, r3's third, won 67
items and lost 36 ($p=0.0015$).\cl{RS28} The plan verdict reads met, 4 of 4 lineages where the rule
needed 3.\cl{RS29} One exam item equalled a program on lineage r2's hand-off list from generation~1, a base
proposal held out in its dev set, not a training program; the examiner excluded it from r2's comparisons and raised its overlap flag on r2's three
verdicts; no other overlap or flag was raised, and no request failed.\cl{RS31,RS40,RS50}

\begin{figure}[tbp]
\centering
\begin{tikzpicture}
\begin{groupplot}[paper, group style={group size=2 by 1, horizontal sep=6mm, y descriptions at=edge left},
  width=0.5\linewidth, height=7.4cm, xmin=-0.06, xmax=0.19, ymin=-14.7, ymax=0.7,
  xtick={-0.05,0,0.05,0.1,0.15}, xticklabel style={/pgf/number format/fixed, /pgf/number format/precision=2},
  ytick={0,-1,-2,-4,-5,-6,-8,-9,-10,-12,-13,-14},
  yticklabels={r0 g1,r0 g2,r0 g3,r1 g1,r1 g2,r1 g3,r2 g1,r2 g2,r2 g3,r3 g1,r3 g2,r3 g3},
  ymajorgrids=false, xlabel={paired gain (share of 889 items)},
  extra x ticks={0}, extra x tick labels={}, extra x tick style={grid style={black!80, thick}},
  error bars/error bar style={line width=0.9pt}, error bars/error mark options={rotate=90, mark size=2.5pt}]
\nextgroupplot[title={S3},
  legend style={at={(1.03,-0.2)}, anchor=north, legend columns=2, /tikz/every even column/.append style={column sep=5mm}}]
\addplot[only marks, mark=*, mark size=2.2pt, okBlue, error bars/.cd, x dir=both, x explicit]
  table[col sep=comma, x=gain, y=y, x error minus=minus, x error plus=plus,
        restrict expr to domain={\thisrow{holds}}{1:1}] {generated/rsi3_forest.csv};
\addplot[only marks, mark=o, mark size=2.4pt, thick, okVermillion, error bars/.cd, x dir=both, x explicit]
  table[col sep=comma, x=gain, y=y, x error minus=minus, x error plus=plus,
        restrict expr to domain={\thisrow{holds}}{0:0}] {generated/rsi3_forest.csv};
\legend{registered sign test holds, registered sign test does not hold}
\nextgroupplot[title={S4}]
\addplot[only marks, mark=*, mark size=2.2pt, okBlue, error bars/.cd, x dir=both, x explicit]
  table[col sep=comma, x=gain, y=y, x error minus=minus, x error plus=plus,
        restrict expr to domain={\thisrow{holds}}{1:1}] {generated/rsi4_forest.csv};
\addplot[only marks, mark=o, mark size=2.4pt, thick, okVermillion, error bars/.cd, x dir=both, x explicit]
  table[col sep=comma, x=gain, y=y, x error minus=minus, x error plus=plus,
        restrict expr to domain={\thisrow{holds}}{0:0}] {generated/rsi4_forest.csv};
\end{groupplot}
\end{tikzpicture}
\caption{Every step's paired gain with the examiner's 95\% interval clustered over the exam's 25
program templates, read from the signed verdicts. Filled points are steps whose registered item-level
sign test holds, open points steps whose test does not. Every first and second S4 step excludes zero;
of the third steps only r2's does, and r0's, r1's and r3's reach below zero, so the third generation's
gain is not robust to the clustering of items by template in three lineages.}\label{fig:rsi-forest}
\end{figure}

\subsection{Template level: the clustered intervals}\label{sec:rsi-robust}

The sign test counts the 889 items as independent, but they come from 25 templates. The examiner also
reports, for every step, the mean per-item paired difference $d_i\in\{-1,0,1\}$ with a cluster-robust
standard error over templates, $\sqrt{\tfrac{G}{G-1}\sum_{c}\big(\sum_{i\in c}(d_i-\bar d)\big)^2}\,/\,n$
with $G=25$, and a two-sided Student-$t$ 95\% interval with $G-1=24$ degrees of freedom
(\cref{fig:rsi-forest}). Every S4 step at generations 1 and 2 excludes zero. At generation~3 only r2's
interval does (0.0811, from 0.0267 to 0.1355); r0's, r1's and r3's include zero. The matching one-sided
95\% bounds are $-0.0054$, $-0.0071$, 0.0360 and 0.0003 for r0 to r3. At the template level, then, the
third generation improves in one lineage (two-sided) or two (one-sided) of four, and a 3-of-4 rule would
not be met; in S3 no lineage clears all three steps.\cl{RS38} The average generation-3 gain over the
lineages is 5.7 points; its clustered interval needs the item-level grades, which the verdicts do not
carry, as they carry no per-template counts.\cl{RS53}

\subsection{What the learner's records show}\label{sec:rsi-process}

\paragraph{Step 1 is inflated by selection.} The exam was selected from items the base got wrong, so
at step 1 the base can lose at most the 12 items it answers, and lost 2, 2, 3 and 1; steps 2 and 3 are
not so constrained.\cl{RS28} The fresh exam, whose items are not selected by the base's answers,
removes this.\cl{RS39}

\paragraph{Churn.} Items right are not cumulative mastery. At step 3 the new adapter lost 28\%, 30\%,
14\% and 19\% of the items its predecessor answered (r0 to r3) while winning 12.4\%, 14.0\%, 13.2\% and
9.5\% of those it missed, although it was trained on a superset of its predecessor's data from the same
base.\cl{RS48} Without a retrain of the same data under another training seed, which is registered
(\cref{sec:followups}), this churn has no noise floor.

\begin{figure}[tbp]
\centering
\begin{tikzpicture}
\begin{axis}[paper, width=0.6\linewidth, height=5.2cm, xmin=-0.25, xmax=3.25, ymin=35, ymax=80,
  xtick={0,1,2,3}, xticklabels={base,1,2,3}, xlabel={generation}, ylabel={dev-set accuracy (\%)},
  legend pos=outer north east, legend cell align=left]
\addplot[lineage r0] table[col sep=comma,x=gen,y=r0] {generated/rsi4_dev.csv};
\addplot[lineage r1] table[col sep=comma,x=gen,y=r1] {generated/rsi4_dev.csv};
\addplot[lineage r2] table[col sep=comma,x=gen,y=r2] {generated/rsi4_dev.csv};
\addplot[lineage r3] table[col sep=comma,x=gen,y=r3] {generated/rsi4_dev.csv};
\legend{r0,r1,r2,r3}
\end{axis}
\end{tikzpicture}
\caption{S4 accuracy on each lineage's own 400-program dev set (100 generator programs and 300 base
proposals, not selected by failure), the absolute effect size the learner's records give.}\label{fig:dev}
\end{figure}

\paragraph{Dev accuracy and frontier yield.} On each lineage's own 400-program dev set, which was not
selected by failure, accuracy rose from 0.4025, 0.40, 0.435 and 0.545 at the base to 0.67, 0.6525, 0.65
and 0.725 at generation~3 (\cref{fig:dev}).\cl{RS46} The champions kept proposing programs they could
not predict: the share of predicted proposals on the frontier fell from 84--93\% at generation~1 to
61--71\% at generation~2 and 47--54\% at generation~3, so even the third champion missed about half of
its own proposals.\cl{RS47}

\paragraph{The dev gate filtered nothing.} As run, the learner's gate advanced every candidate with a
positive net dev gain; significance at its spent alpha (0.025, 0.0083 and 0.0042 at generations 1 to 3)
was recorded, not required. Three of S4's twelve accepted steps (r1's third, r3's second and third) and
six of S3's twelve were not significant at that alpha, and in S3 two accepted steps did not hold on the
exam. At r1's third step the gate saw 25 dev wins against 19 losses ($p=0.226$) while the exam
showed 99 against 55.\cl{RS49,RS24,RS33} TRUST replaces such a greedy gate with an online multiple-testing procedure that
controls the rate of accepted inferior updates;~\citep{trust2026}\cl{CT15} a gate of that kind exists in
the learner as an option but no series used it.\cl{RS34}

\subsection{Registered follow-ups (not yet run)}\label{sec:followups}

Two follow-ups were registered after S4's outcome and before any of their runs, and amended once, before
any verdict.\cl{RS41,RS52} \emph{F}, a re-examination of the base and the twelve S4 champions on an exam
drawn afresh by the examiner (at least 50 templates and 1500 items, not selected by the base's answers),
judged by the registered per-lineage sign test and, co-primary, by the one-sided template-clustered
interval, with the lineage-averaged step gain, strata by old and new template and by the base's answer,
per-template counts, a similarity audit, and descriptive baselines (the base with thinking on, its
four-sample majority answer, and a 27B sibling). \emph{M}, the matched-data control: four lineages in
which the \emph{base} proposes until it holds 2100 frontier programs, trained once from the base, compared
with the four generation-3 champions on F by the difference of mean accuracy with an interval resampling
templates and lineages jointly and an equivalence margin of 2 points, with arms at 900 records (M-900)
and a seed-null retrain of each generation-3 training set.\cl{RS52} Their outcome decides the framing:
generation~3 better than M means the improvement compounds; equivalence means verified data volume
suffices; M better means one round is enough; an inconclusive result leaves the question open. None of
their results is part of this version.\cl{RS39}

\paragraph{Scope.} The result holds for one task family, one base model, three rounds and LoRA training
on interpreter labels, on an exam of 25 templates reused across the four series; it makes no claim of
transfer to other tasks or models.
